\documentclass[12pt]{article} % Times New Roman, 12pt

% SINGLE OR DOUBLE SPACE
%\usepackage{gscale_thesis_singlespace} % Single spaced thesis
\usepackage{../settings/template/gscale_thesis_doublespace} % Double spaced thesis


\usepackage{../settings/template/fancyheadings} % Header and footer styling

\usepackage[numbers, sort&compress]{natbib} % Bibliography

\usepackage{setspace} % Allows double spacing but skips headers/footers
\setcounter{tocdepth}{1} % Limits the TOC to chapter and section names

% Additional packages
\usepackage{graphicx} % Allows the inclusion of figures
\usepackage{subcaption} % Allows captions to be added to subfigures
\usepackage[justification=centering]{caption} % Centres caption text
\usepackage{array} % Used for table formatting
\newcolumntype{P}[1]{>{\raggedright\let\newline\\\arraybackslash\hspace{0pt}}m{#1}}
\usepackage{booktabs} % Fancy-style tables
\usepackage{longtable} % Allows for tables that are more than one page long
\usepackage{float} % Better figure placement control
\usepackage{rotating}
\usepackage{enumerate} % Numbered lists
\usepackage[shortlabels]{enumitem} % For controlling enumerate labels
\usepackage[shortcuts]{extdash} % Allows manual hyphenation of hypenated words
\usepackage{amsmath} % Non-standard math symbols
\usepackage{amsfonts} % Extended fonts for mathematics
\usepackage{multirow}

\usepackage[hidelinks]{hyperref} % Linking to LaTeX labels and external URLs

\numberwithin{equation}{section} % Numbers equations based on their section

\usepackage{tabularx}
\usepackage{pdflscape}
\usepackage{booktabs} % for better horizontal lines
\usepackage{makecell} % for line breaks inside cells

% CUSTOM SETTINGS ********************************
\input{../settings/custom/packages}
\input{../settings/custom/settings}
\newcommand{\indep}{\rotatebox[origin=c]{90}{ $\models$ }}

\title{Monte Carlo estimation for interventional effects with multiple mediators}
\author{Johann Hatzius}
\date{\today}

\begin{document}

\maketitle

\section{Interventional Effects MC Estimation Code}

The \texttt{mc\_mediation\_analysis.R} code implements Monte Carlo simulation based estimators of
three distinct approaches to interventional effects described in the literature for one or two 
mediator problems. The three approaches are that of Lin and Vanderweele (2017) (LV effects), 
Vansteelandt and Daniel (2017) (VD), and Vanderweele and Robinson (2014)/Jackson and Vanderweele
(2018) (VR). Each approach requires different underlying assumptions and targets operationally 
different effects. My understanding of the conceptual machinery underpinning each approach, which
is implemented in my code, is laid out in Section 2. A comparison of the effect estimates from each
approach on how the effect of social class mobility is mediated by IQ and education is given in 
Section 3. I envision integrating Section 2 from this write-up into the section of my dissertation 
that discusses the theory behind mediation analysis with multiple mediators and integrating Section 3
into a section that gives the application of my code and Monte Carlo approach to the social mobility
data. A have produced a table of the arguments for the main MC estimation function I have written
below. 

\begin{center}
\begin{tabular}{ |l|l|p{4cm}| } 
 \hline
 \textbf{Argument} & \textbf{Options} & \textbf{Notes} \\ 
 \hline
 \texttt{data} & dataframe & \\ 
 \hline
 \texttt{outcome\_model} & regression formula & \\ 
 \hline
 \texttt{outcome\_type} & ``continuous'', ``binary'' & \\ 
 \hline
 \texttt{mediator\_models} & vector of formulas & \\ 
 \hline
 \texttt{mediator\_types} & vector of ``continuous''/``binary'' & \\ 
 \hline
 \texttt{te\_model} & regression formula & model for Y(t) - Y(t*)\\ 
 \hline
 \texttt{exposure} & column & \\ 
 \hline
 \texttt{mediators} & vector of columns & \\
 \hline
 \texttt{outcome} & column & \\ 
 \hline
 \texttt{reference} & exposure value & \\
 \hline
 \texttt{effect\_type} & ``LV'', ``VD'', ``VR'' & \\ 
 \hline
 \texttt{vd\_conditional\_model} & regression formula & for VD effects only\\ 
 \hline
 \texttt{estimand} & ``ATE'', ``ATT'', ``ATU'' & ATU should be used for VR efffects \\ 
 \hline
 \texttt{mc\_draws} & integer & \\ 
 \hline
 \texttt{bootstrap\_reps} & integer & \\ 
  \hline
 \texttt{confint} & double (0-1) & \\ 
 \hline

\end{tabular}
\end{center}


\section{Types of Interventional Effects}

\subsection{Lin and Vanderweele (2017)}

The main idea in Lin and Vanderweele (2017) is to estimate path-specific interventional effects 
in the multiple mediator setting in a similar way to how such effects are estimated in the single
mediator case. In the two mediator case, this means that four interventional effects are estimated.
These effects require that the causal ordering of the mediators is known. These effects will sum up 
to the total effect in cases where there are no exposure-mediator or mediator-mediator interactions.
However, these effects will not always sum to the total effect if there are exposure-mediator or 
mediator-interactions. Instead, they always sum to the interventional total effect (defined below). 

Suppose that $M_{1}$ and $M_{2}$ mediate the relationship between an exposure $A$ and an outcome $Y$
and that the causal ordering of $M_{1}$ and $M_{2}$ is known. Suppose $a_{i} = \{0, 1\}$. 
\begin{itemize}
\item Let $M_{1}(a_{1})$ be the distribution of $M_{1}$ for units that receive exposure level $a_{1}$. 
\item Let $M_{2}(a_{1},M_{1}(a_{2}))$ be the distribution of $M_{2}$ for units
that hypothetically receive exposure level $a_{1}$ and a $M_{1}$ level drawn from $M_{1}(a_{2})$.
\item Let $Y(a_{1},M_{1}(a_{2}),M_2(a_{3},M_{1}(a_{4})))$ be the distribution of $M_{2}$ for units
that hypothetically receive exposure level $a_{1}$, a $M_{1}$ level drawn from $M_{1}(a_{2})$,
and an $M_{2}$ level drawn from $M_{2}(a_{3},M_{1}(a_{4}))$.
\end{itemize}

We can write 
$$\psi(a_{1},a_{2},a_{3},a_{4}) = \mathbb{E}(Y|a_{1},a_{2},a_{3},a_{4}) = \int \mathbb{E}(Y|a_{1},m_{1},m_{2})\mathbb{P}(M_{2}|a_{3},m_{1}^{*})\mathbb{P}(M_{1}^{*}|a_{4})\mathbb{P}(M_{1}|a_{2})dM_{1}dM_{1}^{*}dM_{2}$$
where the two versions of $m_{1}$ come from the necessity of separating the $m_{1}$ 
drawn from $M_{1}$ and the $m_{1}$ used to draw from $M_{2}$.   
If we have specified models for the conditional distributions of the mediators and
the outcome, we can use a Monte Carlo approach to estimate $\psi(a_{1},a_{2},a_{3},a_{4})$.
Lin and Vanderweele (2017) then propose the following path-specific effects: 

\begin{itemize}
\item \textbf{Total effect} (TE) --- $Y(1) - Y(0)$: This represents the average causal effect of the change in the exposure 
on changes in the outcome. 
\item \textbf{Total interventional effect} (rTE) --- $\psi(1,1,1,1) - \psi(0,0,0,0)$: This will be the same quantity as the 
total effect if all mediator and outcome models are linear, and will diverge otherwise. The four path-specific effects always
sum to this effect (telescoping). 
\item \textbf{Direct effect} (rDE) --- $\psi(1,0,0,0) - \psi(0,0,0,0)$: This represents the difference between the average outcome
for a unit while intervening to fix two different exposure levels and the mediators at a random level from their joint distribution for
unexposed units.  
\item \textbf{Indirect effect $A \rightarrow M_{1} \rightarrow Y$} (rIE\_M1) --- $\psi(1,1,0,0) - \psi(1,0,0,0)$: This represents 
the difference between the average outcome for a unit between draws from the first mediator distribution given exposure level 1
versus exposure level 0, while holding the exposure fixed at level 1 and drawing the second mediator from its distribution conditional
on levels 0-0 for the exposure and first mediator, respectively. 
\item \textbf{Indirect effect $A \rightarrow M_{2} \rightarrow Y$} (rIE\_M2) --- $\psi(1,1,1,0) - \psi(1,1,0,0)$: This represents 
the difference between the average outcome for a unit between draws from the second mediator distribution given exposure-mediator levels 0-0
versus 1-0, while holding the exposure fixed at level 1 and drawing the first mediator from its distribution conditional
on exposure level 1.
\item \textbf{Indirect effect $A \rightarrow M_{1} \rightarrow M_{2} \rightarrow Y$} (rIE\_M12) --- $\psi(1,1,1,1) - \psi(1,1,1,0)$: This represents 
the difference between the average outcome for a unit between draws from the second mediator distribution given exposure-mediator levels 1-0
versus 1-1, while holding the exposure fixed at level 1 and drawing the first mediator from its distribution conditional
on exposure level 1. 
\end{itemize}


\subsection{Vansteelandt and Daniel (2017)}

The effects in Vansteelandt and Daniel (2017) have the advantage that knowledge of the causal ordering of the mediators is not
required in order to identify them. Additionally, the total effect can be decomposed into the sum of four effects. 

In a similar setup to Lin and Vanderweele (2017), suppose $M_{1}$ and $M_{2}$ mediate the relationship between $A$ and $Y$ but
that the causal ordering of these mediators is unknown. Suppose $a_{i} = \{0, 1\}$. 
\begin{itemize}
\item Let $M_{1}(a_{1})$ be the marginal distribution of $M_{1}$ for units that receive exposure level $a_{1}$. 
\item Let $M_{2}(a_{1})$ be the marginal distribution of $M_{2}$ for units
that receive exposure level $a_{1}$.
\item Let $Y(a_{1},M_{1}(a_{2}),M_2(a_{3}))$ be the distribution of $Y$ for units
that hypothetically receive exposure level $a_{1}$, a $M_{1}$ level drawn from $M_{1}(a_{2})$,
and an $M_{2}$ level drawn from $M_{2}(a_{3})$.
\item Let $Y(a_{1},(M_{1},M_{2})(a_{2}))$ be the distribution of $M_{2}$ for units
that hypothetically receive exposure level $a_{1}$ and a $M_{1}$ level drawn from the joint distribution of $M_{1}, M_{2}$ given
exposure level $a_{2}$. 
\end{itemize}

We can write the formula for the outcome as a conditional distribution over the product of the marginals of the mediators as 
$$\phi(a_{1},a_{2},a_{3}) = \mathbb{E}(Y|a_{1},a_{2},a_{3}) = \int \mathbb{E}(Y|a_{1},m_{1},m_{2})\mathbb{P}(M_{2}|a_{3})\mathbb{P}(M_{1}|a_{2})dM_{1}dM_{2}$$

We can write the formula for the outcome as a conditional distribution over the joint distribution of the mediators as 
$$\tau(a_{1},a_{2}) = \mathbb{E}(Y|a_{1},a_{2}) = \int \mathbb{E}(Y|a_{1},m_{1},m_{2})\mathbb{P}(M_{1}, M_{2}|a_{2})dM_{1}dM_{2}$$

The four effects defined in Vansteelandt and Daniel (2017) are then: 

\begin{itemize}
\item \textbf{Total effect} (TE) --- $Y(1) - Y(0)$: This represents the average causal effect of the change in the exposure 
on changes in the outcome. 
\item \textbf{Direct effect} (rDE) --- $\tau(1,0) - \tau(0,0)$: This represents the difference between the average outcome
for a unit while intervening to fix two different exposure levels and the mediators at a random level from their joint distribution for
unexposed units. Note that in addition to models for the marginals of $M_{1}$ and $M_{2}$ given exposure, a model for either the 
joint distribution of the mediator or for a mediator conditional on the other mediator needs to specified in order to estimate this 
effect. If the same specified conditional model for $M_{2}$ given $M_{1}$ is specified, this will return the same
effect estimate as the direct effect from Lin and Vanderweele (2017). 
\item \textbf{Indirect effect through $M_{1}$} (rIE\_M1) --- $\phi(1,1,0) - \phi(1,0,0)$: This represents 
the difference between the average outcome for a unit between draws from the first mediator distribution given exposure level 1
versus exposure level 0, while holding the exposure fixed at level 1 and drawing the second mediator from its distribution conditional
on level 0 for the exposure.
\item \textbf{Indirect effect through $M_{2}$} (rIE\_M2) --- $\phi(1,0,1) - \phi(1,0,0)$: This represents 
the difference between the average outcome for a unit between draws from the second mediator distribution given exposure levels 0
versus 1, while holding the exposure fixed at level 0 and drawing the first mediator from its distribution conditional
on exposure level 1.
\item \textbf{Remainder}: This represents 
the difference between the total effect and the sum of the direct effect and the two indirect effects. Vansteelandt and Daniel (2017) 
give an interpretation of this quantity as the amount of the effect that is mediated by the dependence between $M_{1}$ and $M_{2}$. 
I am still not sure I fully understand what they are intending, as this effect does not necessarily seem to be related to the
dependence between $M_{1}$ and $M_{2}$. The reason I am unsure on this point is that in the single mediator case, where, using similar
notation to above, the interventional direct effect is given by $\psi(1,0) - \psi(0,0)$ and the indirect effect by $\psi(1,1) - \psi(1,0)$,
it is known that these effects do not necessarily sum up to the total effect (for example, if there are exposure-mediator interactions). 
Then a similar remainder term can be defined, but this would not be interpreted as reflecting dependence between two mediators. Instead,
the right interpretation of such a discrepancy seems to be closer to the interpretation given in Vanderweele (2014), which in the 
context of natural effects in a single mediator problem with sufficient assumptions gives a four-way decomposition of the total effect
in terms of a pure natural direct, a pure indirect effect, an effect of interaction due to presence of the mediator, and an effect of 
interaction not due to presence of the mediator. 
\end{itemize}

\subsection{Vanderweele and Robinson (2014)}

Vanderweele and Robinson (2014) and Jackson and Vanderweele (2018), which builds on the ideas of the first paper in a multiple mediator
situation, give interventional effect estimates for situations in which it does not make substantive sense to consider the exposure as
as hypothetically manipulable quantitiy (e.g. sex, race, potentially social class). In these situations, no decomposition of a total effect
of manipulating the exposure is pursued; instead, a number of interventional effects can be defined that map to the effect on the outcome
for a group with a particular exposure value of equalizing the distribution of that group's mediator value(s) to that of a group with
a different exposure value. The theoretical setup for these effects follows that of Lin and Vanderweele (2017). Unlike for 
the previous two interventional effects, which target the average treatment effect (ATE), this interventional effect targets
the average treatment effect on the untreated (ATU). 
Following Jackson and Vanderweele (2018), the following interventional effects are defined: 

\begin{itemize}
\item \textbf{Total effect} (TE) --- $Y(1) - Y(0)$: This represents the average raw disparity between the outcomes of the 
groups with two different exposure values. It should not be interpreted as a causal effect. 
\item \textbf{Effect of an intervention to equalize $M_{1}$ and keep $M_{2}$ fixed} 
(VRE\_M1\_Marginal) --- $\psi(0,1,0,0) - \psi(0,0,0,0)$: This represents the effect of a hypothetical intervention that would set the 
distribution of $M_{1}$ for a group with exposure level 0 to that of the mediator for a group with exposure 
level 1, while still drawing $M_{2}$ from the distribution of the group with exposure level 0.  
\item \textbf{Effect of an intervention to equalize the marginal of $M_{2}$ and keep $M_{1}$ fixed} 
(VRE\_M2\_Marginal) --- $\psi(0,0,1,1) - \psi(0,0,0,0)$: This represents the effect of a hypothetical intervention that would set the 
distribution of $M_{2}$ for a group with exposure level 0 to that of the mediator for a group with exposure 
level 1, while still drawing $M_{1}$ from the distribution of the group with exposure level 0.  
\item \textbf{Effect of an intervention to equalize $M_{1}$ and draw $M_{2}$ 
from a hypothetical distribution of units with exposure level 0 and $M_{1}$ at exposure level 1}
(VRE\_M1\_Conditional) --- $\psi(0,1,0,1) - \psi(0,0,0,0)$: This represents the effect of a hypothetical intervention that would set the 
distribution of $M_{1}$ for a group with exposure level 0 to that of the mediator for a group with exposure 
level 1, while still drawing $M_{2}$ from the distribution of the group with exposure level 0.  
\item \textbf{Effect of an intervention to equalize the distribution of $M_{2}$ conditional on $M_{1}$ and keep $M_{1}$ fixed} 
(VRE\_M1\_Conditional) --- $\psi(0,0,1,0) - \psi(0,0,0,0)$: This represents the effect of a hypothetical intervention that would set the 
distribution of $M_{2}$ for a group with exposure level 0 to that of the mediator for a hypothetical group with exposure 
level 1 and $M_{1}$ drawn from a distribution with exposure 0,
while still drawing $M_{1}$ from the distribution of the group with exposure level 0.
\item \textbf{Effect of an intervention to equalize the joint distribution of $M_{1}$ and $M_{2}$} --- 
(VRE\_Joint) --- $\psi(0,1,1,1) - \psi(0,0,0,0)$: This represents the effect of a hypothetical intervention that would set the 
joint distribution of the mediators for the group with exposure level 0 to that of the group with exposure level 1. 
\end{itemize}


\section{Application to the Social Mobility Dataset}

I report below the format of my output for the two-mediator interventional effects estimates for the social mobility data for
men in the 1970 cohort,
where IQ and education are the two (potentially causally ordered) mediators of the relationship betwen father's social class
and social class at middle age. Similar to a sensitivity reported in Kuha et al (2021), I create a 5 categories from 
the 7 categories for social class by combining groups 3-5 into one class. I then treat social class as an ordinal, continuous
variable for illustrative purposes. 

I treat all mediators and the outcome as continuous, and I estimate the interventional effects between two exposure levels:
an "untreated" father's social class exposure level of 5 (lowest) and a "treated" father's social class exposure level of 1
(highest). I estimate linear models for each mediator and outcome. 

For each sensitivity I run, I use 1,000,000 Monte Carlo draws and run 500 bootstrap replications. I construct bootstrap
confidence intervals at the $\alpha = 0.05$ confidence level using the percentile bootstrap. 

\subsection{Lin and Vanderweele (2017)}

% latex table generated in R 4.6.0 by xtable 1.8-8 package
% Sat Jun 27 18:56:39 2026
\begin{table}[ht]
\centering
\begin{tabular}{rlrrr}
  \hline
 & Effect & Estimate & 95\% CI LB & 95\% CI UB \\ 
  \hline
1 & TE & -1.14 & -1.28 & -1.01 \\ 
  2 & rTE & -1.14 & -1.28 & -1.01 \\ 
  3 & rDE & -0.56 & -0.69 & -0.42 \\ 
  4 & rIE\_M1 & -0.21 & -0.24 & -0.16 \\ 
  5 & rIE\_M2 & -0.22 & -0.28 & -0.17 \\ 
  6 & r\_IE\_M12 & -0.16 & -0.19 & -0.13 \\ 
   \hline
\end{tabular}
\caption{Linear additive models for mediators, outcome}
\end{table}

% latex table generated in R 4.6.0 by xtable 1.8-8 package
% Sat Jun 27 18:58:16 2026
\begin{table}[ht]
\centering
\begin{tabular}{rlrrr}
  \hline
 & Effect & Estimate & 95\% CI LB & 95\% CI UB \\ 
  \hline
1 & TE & -1.14 & -1.27 & -0.99 \\ 
  2 & rTE & -1.14 & -1.27 & -0.99 \\ 
  3 & rDE & -0.57 & -0.72 & -0.42 \\ 
  4 & rIE\_M1 & -0.19 & -0.27 & -0.12 \\ 
  5 & rIE\_M2 & -0.22 & -0.28 & -0.17 \\ 
  6 & r\_IE\_M12 & -0.16 & -0.18 & -0.13 \\ 
   \hline
\end{tabular}
\caption{Father's class-IQ outcome model interaction term}
\end{table}

% latex table generated in R 4.6.0 by xtable 1.8-8 package
% Sat Jun 27 18:59:55 2026
\begin{table}[ht]
\centering
\begin{tabular}{rlrrr}
  \hline
 & Effect & Estimate & 95\% CI LB & 95\% CI UB \\ 
  \hline
1 & TE & -1.14 & -1.27 & -1.01 \\ 
  2 & rTE & -1.04 & -1.27 & -0.82 \\ 
  3 & rDE & -0.57 & -0.74 & -0.43 \\ 
  4 & rIE\_M1 & -0.19 & -0.27 & -0.12 \\ 
  5 & rIE\_M2 & -0.12 & -0.30 & 0.07 \\ 
  6 & r\_IE\_M12 & -0.16 & -0.19 & -0.13 \\ 
   \hline
\end{tabular}
\caption{Father's class-IQ outcome + education model interaction}
\end{table}

\newpage

\subsection{Vansteelandt and Daniel (2017)}

For both sensitivities, a conditional model for education is specified by regressing education on father's class
and IQ. 

% latex table generated in R 4.6.0 by xtable 1.8-8 package
% Sat Jun 27 19:00:52 2026
\begin{table}[ht]
\centering
\begin{tabular}{rlrrr}
  \hline
 & Effect & Estimate & 95\% CI LB & 95\% CI UB \\ 
  \hline
1 & TE & -1.14 & -1.27 & -1.00 \\ 
  2 & rDE & -0.56 & -0.69 & -0.42 \\ 
  3 & rIE\_M1 & -0.21 & -0.25 & -0.16 \\ 
  4 & rIE\_M2 & -0.38 & -0.44 & -0.31 \\ 
  5 & Remainder & 0.00 & -0.00 & 0.00 \\ 
   \hline
\end{tabular}
\caption{Linear additive models for mediator, outcome}
\end{table}

% latex table generated in R 4.6.0 by xtable 1.8-8 package
% Sat Jun 27 19:05:31 2026
\begin{table}[ht]
\centering
\begin{tabular}{rlrrr}
  \hline
 & Effect & Estimate & 95\% CI LB & 95\% CI UB \\ 
  \hline
1 & TE & -1.14 & -1.29 & -0.98 \\ 
  2 & rDE & -0.56 & -0.69 & -0.42 \\ 
  3 & rIE\_M1 & -0.22 & -0.27 & -0.16 \\ 
  4 & rIE\_M2 & -0.39 & -0.45 & -0.33 \\ 
  5 & Remainder & 0.02 & -0.01 & 0.06 \\ 
   \hline
\end{tabular}
\caption{Father's class-IQ outcome model interaction term}
\end{table}

\subsection{Vanderweele and Robinson (2014)}

% latex table generated in R 4.6.0 by xtable 1.8-8 package
% Sat Jun 27 19:07:39 2026
\begin{table}[ht]
\centering
\begin{tabular}{rlrrr}
  \hline
 & Effect & Estimate & 95\% CI LB & 95\% CI UB \\ 
  \hline
1 & TE & -1.14 & -1.28 & -1.00 \\ 
  2 & rTE & -1.14 & -1.28 & -1.00 \\ 
  3 & VRE\_M1\_Marginal & -0.21 & -0.26 & -0.16 \\ 
  4 & VRE\_M2\_Marginal & -0.38 & -0.44 & -0.32 \\ 
  5 & VRE\_M1\_Conditional & -0.36 & -0.42 & -0.30 \\ 
  6 & VRE\_M2\_Conditional & -0.22 & -0.28 & -0.17 \\ 
  7 & VRE\_Joint & -0.58 & -0.66 & -0.51 \\ 
   \hline
\end{tabular}
\caption{Linear additive models for mediators, outcome}
\end{table}

% latex table generated in R 4.6.0 by xtable 1.8-8 package
% Sat Jun 27 19:08:30 2026
\begin{table}[ht]
\centering
\begin{tabular}{rlrrr}
  \hline
 & Effect & Estimate & 95\% CI LB & 95\% CI UB \\ 
  \hline
1 & TE & -1.14 & -1.27 & -1.00 \\ 
  2 & rTE & -1.14 & -1.27 & -1.00 \\ 
  3 & VRE\_M1\_Marginal & -0.22 & -0.29 & -0.15 \\ 
  4 & VRE\_M2\_Marginal & -0.38 & -0.45 & -0.32 \\ 
  5 & VRE\_M1\_Conditional & -0.38 & -0.46 & -0.30 \\ 
  6 & VRE\_M2\_Conditional & -0.22 & -0.28 & -0.17 \\ 
  7 & VRE\_Joint & -0.60 & -0.70 & -0.50 \\ 
   \hline
\end{tabular}
\caption{Father's class-IQ outcome model interaction term}
\end{table}

% latex table generated in R 4.6.0 by xtable 1.8-8 package
% Sat Jun 27 19:08:54 2026
\begin{table}[ht]
\centering
\begin{tabular}{rlrrr}
  \hline
 & Effect & Estimate & 95\% CI LB & 95\% CI UB \\ 
  \hline
1 & TE & -1.14 & -1.27 & -1.01 \\ 
  2 & rTE & -1.14 & -1.27 & -1.01 \\ 
  3 & VRE\_M1\_Marginal & -0.22 & -0.29 & -0.15 \\ 
  4 & VRE\_M2\_Marginal & -0.38 & -0.44 & -0.31 \\ 
  5 & VRE\_M1\_Conditional & -0.35 & -0.43 & -0.26 \\ 
  6 & VRE\_M2\_Conditional & -0.19 & -0.24 & -0.13 \\ 
  7 & VRE\_Joint & -0.60 & -0.69 & -0.50 \\ 
   \hline
\end{tabular}
\caption{Father's class-IQ outcome + education model interaction}
\end{table}


\end{document}

